% (c) Nikita Lisitsa, lisyarus@gmail.com, 2026

\documentclass[10pt]{beamer}

\usepackage[T2A]{fontenc}
\usepackage[russian]{babel}
\usepackage{minted}

\usepackage[most]{tcolorbox}
\tcbuselibrary{minted}

\usepackage{graphicx}
\graphicspath{ {./images/} }

\usepackage{adjustbox}

\usepackage{color}
\usepackage{soul}
\usepackage{multicol}
\usepackage{hyperref}
\usepackage{tikz}
\usetikzlibrary{arrows.meta,positioning,calc}

\usetheme{metropolis}

\definecolor{red}{rgb}{1,0,0}
\definecolor{codebg}{RGB}{29,35,49}
\setminted{fontsize=\footnotesize}

\makeatletter
\newcommand{\slideimage}[1]{
  \begin{figure}
    \begin{adjustbox}{width=\textwidth, totalheight=\textheight-2\baselineskip-2\baselineskip,keepaspectratio}
      \includegraphics{#1}
    \end{adjustbox}
  \end{figure}
}
\makeatother

\title{Язык программирования C++}
\subtitle{Лекция 15: статические методы, паттерны проектирования, синглтоны, зависимости, автотестирование}
\date{2026}
\author{Никита Лисица}

\setbeamertemplate{footline}[frame number]

\begin{document}

\frame{\titlepage}

\begin{frame}[fragile]
\frametitle{Inline методы класса}
\begin{itemize}
\usemintedstyle{tango}
\item Напоминание: метод, определённый (реализованный) прямо в определении класса, по умолчанию является \mintinline{cpp}|inline|
\usemintedstyle{lightbulb}
\begin{tcblisting}{listing engine=minted, minted language=cpp, minted options={fontsize=\scriptsize}, listing only, colback=codebg, colframe=codebg, rounded corners}
class Array {
private:
  Object* data_;
  size_t size_;

public:
  // Автоматически inline
  size_t size() const {
    return size_;
  }
}
\end{tcblisting}
\end{itemize}
\end{frame}

\begin{frame}[fragile]
\frametitle{Inline методы класса}
\begin{itemize}
\usemintedstyle{tango}
\item Метод, определённый снаружи класса, тоже можно сделать \mintinline{cpp}|inline|:
\usemintedstyle{lightbulb}
\begin{tcblisting}{listing engine=minted, minted language=cpp, minted options={fontsize=\scriptsize}, listing only, colback=codebg, colframe=codebg, rounded corners}
class Array {
private:
  Object* data_;
  size_t size_;

public:
  // Автоматически inline
  size_t size() const;
}

inline size_t Array::size() const {
  return size_;
}
\end{tcblisting}
\end{itemize}
\end{frame}

\begin{frame}[fragile]
\frametitle{Статические переменные}
\begin{itemize}
\usemintedstyle{tango}
\item Напоминание: статическая переменная внутри функции `живёт' (т.е. сохраняет значение, у неё не вызывается деструктор) и между вызовами функции:
\usemintedstyle{lightbulb}
\begin{tcblisting}{listing engine=minted, minted language=cpp, minted options={fontsize=\scriptsize}, listing only, colback=codebg, colframe=codebg, rounded corners}
int gen_unique_id() {
  static int next_id = 0;
  return next_id++;
}
\end{tcblisting}
\pause
\item В каком-то смысле, такая переменная одна на все вызовы функции (т.е. не создаётся каждый раз)
\end{itemize}
\end{frame}

\begin{frame}[fragile]
\frametitle{Статические поля класса}
\begin{itemize}
\usemintedstyle{tango}
\item Поле класса тоже может быть статическим -- тогда оно одно на все объекты (экземпляры) класса:
\usemintedstyle{lightbulb}
\begin{tcblisting}{listing engine=minted, minted language=cpp, minted options={fontsize=\scriptsize}, listing only, colback=codebg, colframe=codebg, rounded corners}
// object.h

class Object {
private:
  // Объявление статического поля
  static int next_id;
  int id;

public:
  Object()
    : id(next_id++)
  {}
};
\end{tcblisting}
\end{itemize}
\end{frame}

\begin{frame}[fragile]
\frametitle{Статические поля класса}
\begin{itemize}
\usemintedstyle{tango}
\item Статическое поле является, по сути, глобальной переменной, привязанной к классу
\pause
\item \(\Longrightarrow\) Под неё нужно выделить место в каком-нибудь cpp-файле:
\usemintedstyle{lightbulb}
\begin{tcblisting}{listing engine=minted, minted language=cpp, minted options={fontsize=\scriptsize}, listing only, colback=codebg, colframe=codebg, rounded corners}
// object.cpp

// Определение статического поля
int Object::next_id = 0;
\end{tcblisting}
\pause
\item Использование:
\usemintedstyle{lightbulb}
\begin{tcblisting}{listing engine=minted, minted language=cpp, minted options={fontsize=\scriptsize}, listing only, colback=codebg, colframe=codebg, rounded corners}
// main.cpp
#include <object.h>

int main() {
  Object o1; // o1.id == 0
  Object o2; // o2.id == 1
}
\end{tcblisting}
\end{itemize}
\end{frame}

\begin{frame}[fragile]
\frametitle{Статические методы класса}
\begin{itemize}
\usemintedstyle{tango}
\item Метод класса тоже может быть статческим
\pause
\item Он тоже `один на все объекты класса`: у него нет \mintinline{cpp}|this|, и он не имеет доступа к нестатическим полям класса:
\usemintedstyle{lightbulb}
\begin{tcblisting}{listing engine=minted, minted language=cpp, minted options={fontsize=\scriptsize}, listing only, colback=codebg, colframe=codebg, rounded corners}
class Object {
private:
  static size_t count_;
  const char* name_;

public:
  Object(const char* name)
    : name_(name)
  {
    ++count_;
  }

  static size_t count() {
    // Нет доступа к this или name_
    return count_;
  }
};
\end{tcblisting}
\end{itemize}
\end{frame}

\begin{frame}[fragile]
\frametitle{Статические методы класса}
\begin{itemize}
\usemintedstyle{tango}
\item Статические методы можно вызывать как через объект класса, так и через сам класс:
\usemintedstyle{lightbulb}
\begin{tcblisting}{listing engine=minted, minted language=cpp, minted options={fontsize=\scriptsize}, listing only, colback=codebg, colframe=codebg, rounded corners}
int main() {
  Object o1;

  // Эти два вызова эквивалентны,
  // есл метод статический:
  o1.count();
  Object::count();
}
\end{tcblisting}
\end{itemize}
\end{frame}

\begin{frame}[fragile]
\frametitle{Паттерны (шаблоны) проектирования}
\begin{itemize}
\usemintedstyle{tango}
\item Повторяющиеся и переиспользуемые архитектурные конструкции, решающие определённые архитектурные задачи
\pause
\item Популярные книги по теме:
\begin{itemize}
\item E. Gamma, R. Helm, R. Johnson, J. Vlissides -- \textit{Design Patterns: Elements of Reusable Object-Oriented Software}
\item R. Martin -- \textit{Clean Code: A Handbook of Agile Software Craftsmanship}
\item R. Nystrom -- \textit{Game Programming Patterns}
\end{itemize}
\pause
\item Все эти паттерны -- \textit{советы, а не строгие правила}, и не стоит пытаться реализовать всё через них
\pause
\begin{tcblisting}{listing engine=minted, minted language=cpp, minted options={fontsize=\scriptsize}, listing only, colback=codebg, colframe=codebg, rounded corners}
R"(
  Sometimes, the elegant implementation is just a function. 
  Not a method. Not a class. Not a framework.
  Just a function.

    - John Carmack
)"
\end{tcblisting}
\end{itemize}
\end{frame}

\begin{frame}[fragile]
\frametitle{Паттерн синглтон (singleton)}
\begin{itemize}
\usemintedstyle{tango}
\item Хотим, чтобы в программе гарантированно был только один объект конкретного класса
\pause
\item Попытка создать два объекта должна быть невозможна (ошибка компиляции)
\pause
\item Примеры:
\pause
\begin{itemize}
\item Класс глобальных настроек
\pause
\item Система логирования
\pause
\item Общий ресурс, используемый разными частями программы
\end{itemize}
\end{itemize}
\end{frame}

\begin{frame}[fragile]
\frametitle{Паттерн синглтон (singleton)}
\begin{itemize}
\usemintedstyle{tango}
\item Нельзя позволять пользователю самому создавать объект \(\Rightarrow\) конструктор должен быть приватным, копирование нужно запретить
\pause
\item Сам единственный экземпляр нужно где-то хранить \(\Rightarrow\) статическая переменная / статическое поле
\pause
\item Нужен доступ к этому экземпляру снаружи \(\Rightarrow\) статические метод, возвращающий ссылку на экземпляр
\end{itemize}
\end{frame}

\begin{frame}[fragile]
\frametitle{Паттерн синглтон (singleton)}
\begin{itemize}
\usemintedstyle{tango}
\item Классический \textit{синглтон Майерса}:
\usemintedstyle{lightbulb}
\begin{tcblisting}{listing engine=minted, minted language=cpp, minted options={fontsize=\scriptsize}, listing only, colback=codebg, colframe=codebg, rounded corners}
class Settings {
private:
  // Приватный конструктор
  Settings();

public:
  // Запрет копирования
  Settings(const Settings&) = delete;
  Settings& operator=(const Settings&) = delete;

  static Settings& instance() {
    static Settings settings;
    return settings;
  }
};
\end{tcblisting}
\end{itemize}
\end{frame}

\begin{frame}[fragile]
\frametitle{Паттерн синглтон (singleton)}
\begin{itemize}
\usemintedstyle{tango}
\item Классический \textit{синглтон Майерса}:
\usemintedstyle{lightbulb}
\begin{tcblisting}{listing engine=minted, minted language=cpp, minted options={fontsize=\scriptsize}, listing only, colback=codebg, colframe=codebg, rounded corners}
int main() {
  Settings::instance().loadFromFile(...);
  ...
}
\end{tcblisting}
\end{itemize}
\end{frame}

\begin{frame}[fragile]
\frametitle{Шаблонный синглтон Майерса}
\begin{itemize}
\usemintedstyle{tango}
\item Синглтонов может быть много, не хочется писать один и тот же код много раз
\pause
\item Сделаем шаблонный класс:
\usemintedstyle{lightbulb}
\begin{tcblisting}{listing engine=minted, minted language=cpp, minted options={fontsize=\scriptsize}, listing only, colback=codebg, colframe=codebg, rounded corners}
template <typename T>
class Singleton {
private:
  Singleton();

public:
  Singleton(const Singleton&) = delete;
  Singleton& operator=(const Singleton&) = delete;

  static T& instance() {
    static T singleton;
    return singleton;
  }
};
\end{tcblisting}
\end{itemize}
\end{frame}

\begin{frame}[fragile]
\frametitle{Шаблонный синглтон Майерса}
\begin{itemize}
\usemintedstyle{lightbulb}
\begin{tcblisting}{listing engine=minted, minted language=cpp, minted options={fontsize=\scriptsize}, listing only, colback=codebg, colframe=codebg, rounded corners}
class Settings {
private:
  ...
public:
  void loadFromFile(...) {...}
};

int main() {
  Singleton<Settings>::instance().loadFromFile(...);
}
\end{tcblisting}
\pause
\usemintedstyle{tango}
\item Класс \mintinline{cpp}|Settings| опять можно создавать и копировать самому
\pause
\item \(\Longrightarrow\) Отнаследуем \mintinline{cpp}|Settings| от \mintinline{cpp}|Singleton<Settings>|
\end{itemize}
\end{frame}

\begin{frame}[fragile]
\frametitle{Curiously Recurring Template Pattern (CRTP)}
\begin{itemize}
\usemintedstyle{lightbulb}
\begin{tcblisting}{listing engine=minted, minted language=cpp, minted options={fontsize=\scriptsize}, listing only, colback=codebg, colframe=codebg, rounded corners}
class Settings : public Singleton<Settings> {
private:
  friend class Singleton<Settings>; 
  Settings() = default;
public:
  void loadFromFile(...) {...}
};
\end{tcblisting}
\pause
\usemintedstyle{tango}
\item Класс \mintinline{cpp}|Settings| нельзя создать самому, только использовать через \mintinline{cpp}|Singleton<Settings>|
\pause
\item Приём наследования класса \mintinline{cpp}|class A : public Base<A>| от шаблонного класса, у которого в качестве шаблонного аргумента передан тот же класс, называется CRTP -- \textit{Curiously Recurring Template Pattern}
\end{itemize}
\end{frame}

\begin{frame}[fragile]
\frametitle{Зависимости кода}
\begin{itemize}
\usemintedstyle{tango}
\item Пусть, в программе есть две функции: \mintinline{cpp}|f()| и \mintinline{cpp}|g()|
\pause
\item Если \mintinline{cpp}|f()| внутри вызывает \mintinline{cpp}|g()|, говорят, что \mintinline{cpp}|f()| \textit{зависит от} \mintinline{cpp}|g()|
\usemintedstyle{lightbulb}
\begin{tcblisting}{listing engine=minted, minted language=cpp, minted options={fontsize=\scriptsize}, listing only, colback=codebg, colframe=codebg, rounded corners}
void f() {
  ...
  g();
  ...
}
\end{tcblisting}
\end{itemize}
\end{frame}

\begin{frame}[fragile]
\frametitle{Зависимости кода}
\begin{itemize}
\usemintedstyle{tango}
\item Аналогично, если в программе есть два класса: \mintinline{cpp}|A| и \mintinline{cpp}|B|
\pause
\item Если класс \mintinline{cpp}|A| использует/содержит класс \mintinline{cpp}|B|, говорят, что \mintinline{cpp}|A| \textit{зависит от} \mintinline{cpp}|B|
\usemintedstyle{lightbulb}
\begin{tcblisting}{listing engine=minted, minted language=cpp, minted options={fontsize=\scriptsize}, listing only, colback=codebg, colframe=codebg, rounded corners}
class A {
private:
  B* b;
public:
  ...
};
\end{tcblisting}
\end{itemize}
\end{frame}

\begin{frame}[fragile]
\frametitle{Зависимости кода}
\begin{itemize}
\usemintedstyle{tango}
\item Почему полезно думать про зависимости кода?
\pause
\item Если мы хотим в новую библиотеку/программу/подсистему добавить некую функцию/класс, то нужно добавить и все его зависимости
\pause
\item \(\Longrightarrow\) Хочется, чтобы зависимостей было меньше
\pause
\item При этом, в зависимостях часто лежит переиспользуемый, обобщённый код
\pause
\item \(\Longrightarrow\) Хочется, чтобы зависимостей было больше :)
\end{itemize}
\end{frame}

\begin{frame}[fragile]
\frametitle{Циклические зависимости}
\begin{itemize}
\usemintedstyle{tango}
\item Если две части программы (функции/классы/файлы/etc) зависят друг от друга, говорят о \textit{циклической зависимости}
\pause
\item Это почти всегда плохо: эти части нельзя использовать по отдельности (независимо), возникают проблемы с порядком определения, и т.п.
\usemintedstyle{lightbulb}
\begin{tcblisting}{listing engine=minted, minted language=cpp, minted options={fontsize=\scriptsize}, listing only, colback=codebg, colframe=codebg, rounded corners}
// Приходится делать forward declaration,
// иначе поле A::b имеет неизвестный тип
class B;

class A {
private:
  B* b;
};

class B {
private:
  A* a;
};
\end{tcblisting}
\end{itemize}
\end{frame}

\begin{frame}[fragile]
\frametitle{Циклические зависимости}
\begin{itemize}
\usemintedstyle{tango}
\item Часто циклическую зависимость можно устранить
\pause
\item Объединить созависимые компоненты в одну

\begin{center}
\begin{tikzpicture}[node distance=2cm, >=Stealth]
  \node (A) {$A$};
  \node (B) [right=of A] {$B$};
  \node (C) [right=of B] {$A+B$};
  \draw[->] (A) to[bend left=20] (B);
  \draw[->] (B) to[bend left=20] (A);
  \draw[->, double, shorten >=0.5cm, shorten <=0.5cm] (B) to (C);
\end{tikzpicture}
\end{center}

\pause
\item Выделить общую часть, от которой будут зависеть эти компоненты

\begin{center}
\begin{tikzpicture}[node distance=2cm, >=Stealth]
  \node (A) {$A$};
  \node (B) [right=of A] {$B$};
  \node (C) [right=of B] {$$};
  \node (D) [below=0.5cm of C] {$A$};
  \node (E) [right=of D] {$B$};
  \node (F) [above=1cm of $(D)!0.5!(E)$] {$C$};
  \draw[->] (A) to[bend left=20] (B);
  \draw[->] (B) to[bend left=20] (A);
  \draw[->, double, shorten >=0.5cm, shorten <=0.5cm] (B) to (C);
  \draw[->] (F) to (D);
  \draw[->] (F) to (E);
\end{tikzpicture}
\end{center}
\end{itemize}
\end{frame}

\begin{frame}[fragile]
\frametitle{Пример архитектуры программы: Гомоку}
\begin{itemize}
\usemintedstyle{tango}
\item Задача: спроектировать и реализовать игру \textit{гомоку} (как крестики-нолики, 5 в ряд, на доске 15\(\times\)15)
\pause
\item Спроектировать = продумать файлы/функции/классы/модули и их взаимодействие
\pause
\item Избежать излишних зависимостей (особенно циклических) между компонентами программы
\pause
\item Реализовать автоматическое тестирование (т.н. \textit{автотесты})
\end{itemize}
\end{frame}

\begin{frame}[fragile]
\frametitle{Архитектура программы: паттерн Model-View}
\begin{itemize}
\usemintedstyle{tango}
\item Компонента \textit{Model}
\begin{itemize}
\item Хранение внутреннего представления данных программы
\item Контроль правил игры и реализация ходов игроков
\item Обеспечение инвариантов внутреннего состояния
\end{itemize}
\pause
\item Компонента \textit{View}
\begin{itemize}
\item Взаимодействие с пользователем
\item Ввод ходов пользователя
\item Вывод состояния на экран
\end{itemize}
\pause
\item Кто от кого должен зависеть? \pause Подумаем о повторном использовании кода
\end{itemize}
\end{frame}

\begin{frame}[fragile]
\frametitle{Архитектура программы: паттерн Model-View}
\begin{itemize}
\usemintedstyle{tango}
\item Переиспользование \textit{Model} -- напишем разные \textit{View}, использующие один \textit{Model}:
\begin{itemize}
\item Текстовый/графический интерфейс под разные платформы
\item Веб-интерфейс
\item Мобильное приложение
\end{itemize}
\pause
\item Переиспользование \textit{View} -- напишем универсальный \textit{View}, работающий с разными моделями игр \textit{Model}:
\begin{itemize}
\item Гомоку
\item Шашки
\item Шахматы
\item Го
\end{itemize}
\pause
\item Скорее всего, разные игры требуют и разного ввода/вывода, так что вариант с переиспользованием \textit{View} не реалистичен
\pause
\item \(\Longrightarrow\) Сделаем переиспользуемый \textit{Model}, и от него будет зависеть \textit{View}
\end{itemize}
\end{frame}

\begin{frame}[fragile]
\frametitle{Другие Model-View паттерны}
\begin{itemize}
\usemintedstyle{tango}
\item Model-View-Controller (MVC)
\item Model-View-Presenter (MVP)
\item Model-View-ViewModel (MVVM)
\end{itemize}
\end{frame}

\begin{frame}[fragile]
\frametitle{Гомоку: Model}
\begin{itemize}
\usemintedstyle{tango}
\item Что входит в состояние игры?
\pause
\begin{itemize}
\item Чей сейчас ход / чем закончилась игра
\pause
\item Состояние всех клеток
\end{itemize}
\pause
\usemintedstyle{lightbulb}
\begin{tcblisting}{listing engine=minted, minted language=cpp, minted options={fontsize=\scriptsize}, listing only, colback=codebg, colframe=codebg, rounded corners}
enum class GameState {
  XPlayerTurn, // ход крестиков
  OPlayerTurn, // ход ноликов
  XPlayerWin,  // крестики выиграли
  OPlayerWin,  // нолики выиграли
  Draw,        // ничья
};

// 'O' или 'X'
using Player = char;
\end{tcblisting}
\end{itemize}
\end{frame}

\begin{frame}[fragile]
\frametitle{Гомоку: Model}
\begin{itemize}
\usemintedstyle{tango}
\usemintedstyle{lightbulb}
\begin{tcblisting}{listing engine=minted, minted language=cpp, minted options={fontsize=\scriptsize}, listing only, colback=codebg, colframe=codebg, rounded corners}
Player currentPlayer(GameState state) {
  switch (state) {
    case GameState::XPlayerTurn: return 'X';
    case GameState::OPlayerTurn: return 'O';
    default: return 0;
  }
}

GameState switchPlayer(GameState state) {
  switch (state) {
    case GameState::XPlayerTurn: return GameState::OPlayerTurn;
    case GameState::OPlayerTurn: return GameState::XPlayerTurn;
    default: state;
  }
}
\end{tcblisting}
\end{itemize}
\end{frame}

\begin{frame}[fragile]
\frametitle{Гомоку: Model}
\begin{itemize}
\usemintedstyle{lightbulb}
\begin{tcblisting}{listing engine=minted, minted language=cpp, minted options={fontsize=\scriptsize}, listing only, colback=codebg, colframe=codebg, rounded corners}
class Model {
private:
  char board_[15][15];
  GameState state_ = GameState::XPlayerTurn;

  // После хода по координатам (x, y) проверить,
  // не окончена ли игра
  void checkWinner(int x, int y);

public:
  // Проинициализировать пустое поле
  Model();

  // Сделать ход и обновить состояние,
  // либо сообщить, что ход неверный
  bool move(int x, int y);
  char get(int x, int y) const;

  // Проверить, закончена ли игра
  bool finished() const;
  // Сообщить, есть ли победитель
  std::optional<Player> winner() const;
};
\end{tcblisting}
\end{itemize}
\end{frame}

\begin{frame}[fragile]
\frametitle{Гомоку: Model}
\begin{itemize}
\usemintedstyle{lightbulb}
\begin{tcblisting}{listing engine=minted, minted language=cpp, minted options={fontsize=\scriptsize}, listing only, colback=codebg, colframe=codebg, rounded corners}
bool Model::move(int x, int y) {
  // Игра уже окончена
  if (finished())
    return false;

  // Клетка уже заполнена
  if (board_[x][y] != ' ')
    return false;

  board_[x][y] = currentPlayer(state_);

  checkWinner(x, y);

  if (!finished())
    state_ = switchPlayer(state_);
}
\end{tcblisting}
\end{itemize}
\end{frame}

\begin{frame}[fragile]
\frametitle{Гомоку: View}
\begin{itemize}
\usemintedstyle{lightbulb}
\begin{tcblisting}{listing engine=minted, minted language=cpp, minted options={fontsize=\scriptsize}, listing only, colback=codebg, colframe=codebg, rounded corners}
class View {
private:
  Model& model_;

public:
  View(Model& model);

  void processInput(std::istream& in, std::ostream& out);
  void printBoard(std::ostream& out);
  void showWinner(std::ostream& out);
};
\end{tcblisting}
\end{itemize}
\end{frame}

\begin{frame}[fragile]
\frametitle{Гомоку: View}
\begin{itemize}
\usemintedstyle{lightbulb}
\begin{tcblisting}{listing engine=minted, minted language=cpp, minted options={fontsize=\scriptsize}, listing only, colback=codebg, colframe=codebg, rounded corners}
void View::processInput(std::istream& in, std::ostream& out) {
  while (true) {
    int x, y;
    in >> x >> y;
    if (model_.move(x, y))
      break;
    out << "Invalid move" << std::endl;
  }
}

void View::printBoard(std::ostream& out) {
  for (int y = 0; y < 15; ++y) {
    for (int x = 0; x < 15; ++x) {
      out << model_.get(x, y);
    }
    out << '\n';
  }
  out << std::flush;
}
\end{tcblisting}
\end{itemize}
\end{frame}

\begin{frame}[fragile]
\frametitle{Гомоку: main}
\begin{itemize}
\usemintedstyle{lightbulb}
\begin{tcblisting}{listing engine=minted, minted language=cpp, minted options={fontsize=\scriptsize}, listing only, colback=codebg, colframe=codebg, rounded corners}
int main() {
  Model model;
  View view(model);
  while (!model.finished()) {
    view.printBoard();
    view.processInput(std::cout);
  }
  view.showWinner(std::cout);
}
\end{tcblisting}
\end{itemize}
\end{frame}

\begin{frame}[fragile]
\frametitle{Тестирование}
\begin{itemize}
\usemintedstyle{tango}
\item Любой программный продукт тестируется специально обученными людьми
\pause
\item Часто, процесс тестирования можно (и нужно!) автоматизировать:
\pause
\begin{itemize}
\item Проверить, что функция на известных входах даёт известные выходы
\pause
\item Проверить, что сохраняются инварианты класса
\pause
\item Проверить, что типовые сценарии приводят к ожидаемому результату
\pause
\item Проверить, что правильно обрабатываются некорректные входные данные
\pause
\item Проверить, что после очередных изменений в коде ничего не сломалось
\end{itemize}
\pause
\item \(\Longrightarrow\) \textit{Автоматические тесты (автотесты)}
\end{itemize}
\end{frame}

\begin{frame}[fragile]
\frametitle{Автотестирование}
\begin{itemize}
\usemintedstyle{tango}
\item Тест -- это отдельная программа, выполняющая тестирование
\pause
\item Часто все сценарии тестирования собраны в одну тестирующую программу
\pause
\item Каждый сценарий -- отдельная функция, проверяющая конкретное поведение
\pause
\item Есть множество библиотек для автотестирования, упрощающих этот процесс: GoogleTest, Boost.Test, Catch2, doctest, QuickCheck
\end{itemize}
\end{frame}

\begin{frame}[fragile]
\frametitle{Автотестирование}
\begin{itemize}
\usemintedstyle{lightbulb}
\begin{tcblisting}{listing engine=minted, minted language=cpp, minted options={fontsize=\scriptsize}, listing only, colback=codebg, colframe=codebg, rounded corners}
bool testSameMove() {
  Model model;
  // Делаем один и тот же ход два раза
  model.move(0, 0);
  // Второй раз move должна вернуть false
  if (model.move(0, 0)) {
    // Тест не прошёл
    return false;
  }
  // Тест прошёл
  return true;
}
\end{tcblisting}
\end{itemize}
\end{frame}

\begin{frame}[fragile]
\frametitle{Автотестирование}
\begin{itemize}
\usemintedstyle{lightbulb}
\begin{tcblisting}{listing engine=minted, minted language=cpp, minted options={fontsize=\scriptsize}, listing only, colback=codebg, colframe=codebg, rounded corners}
bool testHorizontalWin() {
  Model model;
  // Делаем простые ходы
  for (int i = 0; i < 4; ++i) {
    // Ход крестиков
    model.move(i, 0);
    // Ход ноликов
    model.move(i, 1);
  }
  // Победный ход крестиков
  model.move(4, 0);

  // Условие прохождения теста
  return model.finished() && model.winner() == 'X';
}
\end{tcblisting}
\end{itemize}
\end{frame}

\begin{frame}[fragile]
\frametitle{Автотестирование}
\begin{itemize}
\usemintedstyle{lightbulb}
\begin{tcblisting}{listing engine=minted, minted language=cpp, minted options={fontsize=\scriptsize}, listing only, colback=codebg, colframe=codebg, rounded corners}
using TestFunction = bool (*)();

struct TestData {
  TestFunction function;
  const char* name;
};

class TestRunner {
private:
  std::vector<TestData> tests_;

public:
  void add(TestFunction function, const char* name);
  void run();
};

int main() {
  TestRunner runner;
  runner.add(testSameMove, "testSameMove");
  runner.add(testHorizontalWin, "testHorizontalWin");
  runner.run();
}
\end{tcblisting}
\end{itemize}
\end{frame}

\begin{frame}[fragile]
\frametitle{Автотестирование}
\begin{itemize}
\usemintedstyle{lightbulb}
\begin{tcblisting}{listing engine=minted, minted language=cpp, minted options={fontsize=\scriptsize}, listing only, colback=codebg, colframe=codebg, rounded corners}
void TestRunner::run() {
  int success = 0;
  
  for (const TestData& test : tests_) {
    if (test.function()) {
      std::cout << "Test " << test.name << ": success\n";
      ++success;
    } else {
      std::cout << "Test " << test.name << ": failure\n";
    }
  }

  std::cout << "Success: " << success << "/"
    << tests_.size() << "\n";
}
\end{tcblisting}
\end{itemize}
\end{frame}

\end{document}